\chapter*{使用指南}
\addcontentsline{toc}{chapter}{使用指南}
\markboth{使用指南}{使用指南}
\begin{flushright}——写给刑法小白的一本书\end{flushright}
\section*{一、内容说明}
\begin{enumerate}[label=\arabic*.,leftmargin=2.2em]
\item 适用人群：刑法小白，零基础。
\item 讲解特色：注重从“\CJKunderwave{底层原理}”$\rightarrow$“\CJKunderwave{行为模型}”。基于此，才能将“知识点”$\rightarrow$“答题能力”。
\end{enumerate}
\section*{二、通关产品}
\begin{enumerate}[label=\arabic*.,leftmargin=2.2em]
\item 《精讲卷》，系统全面的基础教材，必须认认真真看两遍。
\item 《真金题卷》，这是复习的黄金资料，必须扎扎实实刷两遍。
\item 《背诵卷》，提炼年度精华考点，归纳总结。
\item 《考前 100 题》，精心编制模拟题，实战演练。
\end{enumerate}
\section*{三、交流渠道}
新浪微博：“柏浪涛的刑法观”。

众合在线 APP，竹马 APP。

\chapter*{法考是不是“三无考试”？}
\addcontentsline{toc}{chapter}{法考是不是“三无考试”？}
\markboth{法考是不是“三无考试”？}{法考是不是“三无考试”？}
\section*{一、法考没有官方教材}
如果有，就违法了。这是因为，《行政许可法》第 54 条中规定：“行政机关或者行业组织应当事先公布资格考试的报名条件、报考办法、考试科目以及考试大纲。但是，不得组织强制性的资格考试的考前培训，\CJKunderwave{不得指定教材或者其他助考材料}。”

根据该规定，法考这种职业资格考试实行“管—培分离”制度，也即司法部负责“考什么”，不负责指导考生“学什么”，不负责指导培训讲师“教什么”。“学什么”和“教什么”由考生和讲师自己把握，自己负责。考生备考和讲师辅导是个自由竞争的市场，不存在行政垄断。

\section*{二、九大本教材的地位}
虽然业内（包括我）都习惯称九大本教材是官方教材，但严格讲，九大本教材不是司法部的指定教材，具有参考价值（因为是知名学者编写、出版质量精良），但没有法定效力。

\section*{三、马工程教材的地位}
根据《行政许可法》第 54 条规定，马工程教材不是法考的官方教材，具有参考价值，但没有法考领域的法定效力。该教材对本科学历教育具有指导地位。

\section*{四、法考没有官方公布的真题和答案}
自 2017 年后，司考改革为法考，司法部不再公布真题和答案。

\section*{五、法考培训讲师的处境}
由于\CJKunderwave{没有官方教材}，也\CJKunderwave{没有了官方公布的真题和答案}，也即没有了公开的统一标准，法考成了“\CJKunderwave{三无考试}”。讲师只能自己摸索把握讲解内容和观点，猜测官方观点。任何人都有权去猜测（当然，内容观点不得违反法律和司法解释，这是不言而喻的）。

猜测的效果好不好，考生是唯一有效的评价主体，考生根据成绩进行评价，是唯一有效的评价标准。

由此\CJKunderwave{一方面，讲师享有讲解内容的选择权，讲解的观点不存在违反某个统一标准的问题，另一方面，讲师对自己选择的观点要承担考生用脚投票的后果}。自己选择，自己负责。

讲师授课如同在大海航行，没有官方灯塔导航，全凭自己\CJKunderwave{摸黑前行}。如果偏航了，就自己校正，没必要怨天尤人。如果方向对了，也没必要沾沾自喜，因为很大程度是运气。

在这法考的艰难航行中，考生和讲师都属于被害人自陷风险，自己负责。唯有相扶相携，才能顺利上岸！

\chapter*{知识体系}
\addcontentsline{toc}{chapter}{知识体系}
\markboth{知识体系}{知识体系}
\section*{总论}
\InsertDiagram{figures/knowledge_general.tex}
\section*{分论}
\InsertDiagram{figures/knowledge_specific.tex}
